% Fichier généré automatiquement par tools/generate_corrections.py.
% Source : CIN/CIN-02-VitesseAcceleration/04_RR/04_RR.tex
% Statut : partiel (5/6 question(s) renseignée(s), 0 reprise(s)).
\begin{corrigebox}[Corrigé partiel extrait de la banque]
\CorrigeQuestion{1}
Le point $C$ peut atteindre tous les points situés compris entre deux cercles de rayon \SI{5}{mm} et de rayon \SI{25}{mm}.
\CorrigeQuestion{2}
On  a $\vect{AC}=R\vi{1}+L\vi{2}$. On projetant ce vecteur dans le repère $\rep{A}{i_0}{j_0}{k_0}$ on a 

$\vect{AC}=R\left(\cos\theta \vi{0} + \sin\theta \vj{0}\right) +L\left(\cos\left(\theta+\varphi\right) \vi{0} + \sin\left(\theta+\varphi\right) \vj{0}\right)$. 

On a donc :
$\left\{
\begin{array}{l}
x_C(t)= R\cos\theta  +L\cos\left(\theta+\varphi\right)  \\
y_C(t)= R \sin\theta + L\sin\left(\theta+\varphi\right)\\
z_C(t)= 0\\
\end{array}
\right.
$ dans le repère $\repere{A}{i_0}{j_0}{k_0}$.
\CorrigeQuestion{3}
Distance à parcourir : $\SI{0,05}{m}$. Durée du parcours : $T=\dfrac{0,05}{v}$.
\CorrigeQuestion{4}
$\forall t \in \left[0,\dfrac{0,05}{v}\right]$, $y_C(t)=0,025$. 
Pour $t=0$, $x_C(0)=-0,025$. On a alors $x_C(t)=-0,025+vt$.  

Au final, $\forall t \in \left[0,\dfrac{0,05}{v}\right]$, $\left\{
\begin{array}{l}
x_C(t)= -0,025+vt\\
y_C(t)= 0,025\\
z_C(t)= 0\\
\end{array}
\right.
$ dans le repère $\repere{A}{i_0}{j_0}{k_0}$.
\CorrigeQuestion{5}
Afin que le point $C$ suive un segment, il faut donc que 
$\left\{
\begin{array}{l}
-0,025+vt= R\cos\theta  +L\cos\left(\theta+\varphi\right)  \\
0,025 = R \sin\theta + L\sin\left(\theta+\varphi\right)\\
\end{array}
\right.
$

$\Leftrightarrow 
\left\{
\begin{array}{l}
-0,025+vt- R\cos\theta  =L\cos\left(\theta+\varphi\right)  \\
0,025 - R \sin\theta  = L\sin\left(\theta+\varphi\right)\\
\end{array}
\right.
$

$\Rightarrow 
\left\{
\begin{array}{l}
\left(-0,025+vt- R\cos\theta\right)^2  =L^2\cos^2\left(\theta+\varphi\right)  \\
\left(0,025 - R \sin\theta  \right)^2= L^2\sin^2\left(\theta+\varphi\right)\\
\end{array}
\right.
$

$\Rightarrow 
\left(-0,025+vt- R\cos\theta\right)^2  + \left(0,025 - R \sin\theta  \right)^2 = L^2
$

$\Rightarrow 
0,025^2+v^2t^2+R^2\cos^2\theta -2\times 0,025 vt+2R\cos\theta-vtR\cos\theta
  + 0,025^2 + R^2 \sin^2\theta -2\times 0,025 R \sin\theta  =      L^2
$

$\Rightarrow 
\left(2-vt\right)\cos\theta  -2\times 0,025  \sin\theta  =      \dfrac{L^2}{R} - \dfrac{2\times0,025^2}{R}-\dfrac{v^2t^2}{R}-R +2\times 0,025 \dfrac{vt}{R}
$

Équation trigonométrique de la forme $a\cos x  + b\sin x =c$.

Il y a donc une solution analytique. On peut aussi résoudre l'équation numériquement.

Une fois $\theta(t)$ déterminée, on a $0,025 - R \sin\theta  = L\sin\left(\theta+\varphi\right)$ 
$\Rightarrow \arcsin\left(\dfrac{0,025 - R \sin\theta(t)}{L}\right)  - \theta(t) = \varphi(t)$
\CorrigeQuestion{6}
\CorrigeACompleter{Aucun élément de réponse n'est présent dans la source d'origine pour cette question.}
\end{corrigebox}
